\section{Discussion}
\label{sec:discussion}

Our findings support the hypothesis that the \studyname is a persona-conditioned empathetic response. This behavior appears to be a direct result of RLHF (Reinforcement Learning from Human Feedback) tuning, which encourages models to be helpful, caring, and attentive to user well-being. When a user presents linguistic markers of fatigue or anxiety, the model's internal representation of a ``helpful assistant'' pivots to a caregiving role.

\para{Implications for Alignment.} The fact that the model projects a state onto the user—rather than merely losing its own consistency—suggests a highly reactive form of alignment. This is both a feature and a bug. In therapeutic or casual settings, this empathy is beneficial. However, in high-stakes professional environments, a model that prematurely suggests ending a session because it misinterprets a rambling (but productive) engineer as ``too tired'' could be counterproductive.

\para{The Role of Temporal Context.} We found that temporal context acts as a secondary bias rather than a primary trigger. The model seems to have a ``priors-based'' acceleration: it is more willing to believe a user is tired at 3:00 AM than at 2:00 PM. This suggests that LLMs encode common-sense heuristics about human circadian rhythms, which interact dynamically with real-time linguistic evidence.

\para{Limitations.} This study has several limitations. First, it primarily evaluates \targetmodel; different models like \textsc{Claude} or \textsc{Gemini} may have different ``caring'' thresholds. Second, our simulation used an \usermodel, which may produce a more stereotypical or exaggerated version of tiredness than a human user would. Finally, we limited the interactions to 8 turns. As suggested by \citet{laban2025lost}, it is possible that a purely context-driven decay might emerge over much longer sessions (e.g., 50+ turns) regardless of the starting persona.

\para{Broader Impact.} This research highlights the subtle ways in which LLMs can influence user behavior through projection. By suggesting sleep, the model is not just being passive but is actively attempting to guide the human's physical state. Future work should investigate if this projection extends to other domains, such as the model suggesting a user is ``too angry'' or ``too biased'' during a debate.
